
\section{问题一：质量冲突消解与综合评价模型}
\label{sec:q12}

\subsection{文献依据与模型选择}

多指标评价的冲突可能来自同一属性的测量分歧，也可能来自不同属性之间的取舍。文献的部分研究决策研究采用极小极大方法协调冲突\cite{yu2016}；部分研究强调辨别冲突信息的来源，再确定融合权重\cite{wang2023}；部分研究将冲突测度与决策目标相联系\cite{chen2026}。这些研究对象不同于文本质量，但共同表明冲突分型应先于聚合规则的选择。

第一小问采用的 CRITIC 权重由指标离散度及相关性确定\cite{diakoulaki1995}，然而本题中的“指标冲突”描述的是样本间统计关系而非只关乎权重。并且存在于同一个文本中存在一个评价方向高分而另一个评价方向低分，故采用加权平均来用高分弥补低分，其补偿程度需单独建模\cite{giarlotta2001}。FineWeb 和 DataComp-LM 的实验表明数据过滤方式会影响语言模型表现\cite{penedo2024,li2024}。因而保留原 CRITIC--TOPSIS 分数，另构造冲突标记与有限补偿评价分。

\subsection{数据衔接与冲突定义}

设 $x_{ij}\in[0,1]$ 为方向统一后的第 $j$ 项指标，$w_j\geq0$ 为第一小问权重。把 11 项较通用指标归为五组 $G_g$：内容价值、表达质量、整洁度、无广告程度和非重复程度。组内均分与组权重定义为：
\begin{equation}
z_{ig}=\frac{\sum_{j\in G_g}w_jx_{ij}}{\sum_{j\in G_g}w_j},\qquad
v_g=\frac{\sum_{j\in G_g}w_j}{\sum_{h=1}^{5}\sum_{j\in G_h}w_j},
\qquad \sum_{g=1}^{5}v_g=1.
\label{eq:q12-group}
\end{equation}

内容价值包括 FineWeb 教育价值、QuRating 教育价值及事实与常识代理。表达质量包括英语流畅度、可读性及写作风格；非重复程度包括独特词比例及两项反向高频词组指标；第一小问原评分及全指标对照中仍保留DSIR 三项、无字母词比例和句末标点比例。

在 A1 拟合集计算并选定每组的 $20\%$ 与 $80\%$ 分位阈值 $l_g,h_g$。定义高低集合与冲突标记：
\begin{equation}
H_i=\{g:z_{ig}>h_g\},\qquad
L_i=\{g:z_{ig}<l_g\},\qquad
I_i=\mathbf{1}\{|H_i|>0\ \land\ |L_i|>0\}.
\label{eq:q12-conflict}
\end{equation}

严格不等号使处于阈值上的样本归入中间状态。$I_i=1$ 只表示同一文本至少有一组相对突出、另一组相对偏低。由 $H_i,L_i$ 表示是否为空，可得到冲突、仅优势、仅短板、中间状态四个互斥且穷尽的类别。冲突程度定义为：
\begin{equation}
C_i=\max_g\left[\frac{z_{ig}-h_g}{1-h_g}\right]_+
    \max_g\left[\frac{l_g-z_{ig}}{l_g}\right]_+,
\qquad
B_i=\frac{|H_i||L_i|}{\binom{5}{2}}.
\label{eq:q12-strength}
\end{equation}
其中 $C_i$ 为跨阈幅度，$B_i$ 为高低组配对广度。阈值来自 A1 拟合集，数值见表~\ref{tab:q12-thresholds}。

\begin{table}[htbp]
\centering
\caption{五维冲突判别阈值与权重}
\label{tab:q12-thresholds}
\begin{tabular}{lccc}
\toprule
维度 & 低阈值 $l_g$ & 高阈值 $h_g$ & 组权重 $v_g$ \\
\midrule
内容价值   & 0.2783 & 0.5613 & 0.2499 \\
表达质量   & 0.3578 & 0.7702 & 0.2836 \\
整洁度     & 0.2673 & 0.9884 & 0.1090 \\
无广告程度 & 0.5244 & 0.9004 & 0.1056 \\
非重复程度 & 0.6366 & 0.7838 & 0.2519 \\
\bottomrule
\end{tabular}
\end{table}

\subsection{冲突成因与消解规则}

在选定的评价质量指标冲突情况归纳中存在三种不同类型冲突，且对于不同类型的冲突应采取不同处理。第一种冲突为\textbf{测量冲突}，如两种教育维度对同一属性展现出一高一低。第二种冲突为\textbf{属性权衡冲突}，如内容价值高并且广告信号强、整洁度高而非重复程度低，考虑通过限制不同维度的相互补偿表达决策偏好。第三种冲突\textbf{适用性冲突}，即代码、非英语和 LaTeX 文本的英语流畅度高但是格式分低。

A1 中两种教育评分一高一低的样本为 415 条（$0.8101\%$），全量去重样本中有 1,402 条（$0.5370\%$）；组内分歧可能被五维平均后被掩盖，需要分开分析。以内容价值高且无广告程度低判别，A1 命中 1,884 条；换成原始广告 logit 大于无广告 logit，仅有 310 条。15 条定向原文案例虽然揭示了上述可能机制，但不能估计总体误判率。

\subsection{有限补偿综合评价模型}

第一小问的原分 $Q_i^{(0)}$ 作为参照。为处理属性间权衡采用名义权重来计算，即采用五个维度的各组权重；计算方法采用先把同一组内的 CRITIC 指标权重相加，再除以五组所包含指标的权重总和所得到的权重。先构造五维加权平均 $M_i=\sum_gv_gz_{ig}$，再允许组权重在名义权重 $v$ 附近变化。令 $\Delta_5=\{p:p_g\geq0,\sum_gp_g=1\}$，定义：
\begin{equation}
\mathcal W_{\lambda}=\{(1-\lambda)v+\lambda p:p\in\Delta_5\},\qquad
Q_i^{\mathrm{rob}}(\lambda)=100\min_{a\in\mathcal W_{\lambda}}\sum_{g=1}^{5}a_gz_{ig},
\quad 0\leq\lambda\leq1.
\label{eq:q12-robust}
\end{equation}

$\lambda$ 是对组权重不确定性的保守程度。内层目标关于 $p$ 为线性函数，在单纯形顶点取得最小值，因此若 $m_i=\min_g z_{ig}$，便有闭式解：
\begin{equation}
Q_i^{\mathrm{rob}}(\lambda)=100\big[(1-\lambda)M_i+\lambda m_i\big],\qquad
Q_i^{\mathrm{rob}}(\lambda)-100M_i=-100\lambda(M_i-m_i).
\label{eq:q12-closed}
\end{equation}

$\lambda=0$ 是五维加权平均，$\lambda=1$ 是最弱维度；对任意 $0\leq\lambda\leq1$，有 $100m_i\leq Q_i^{\mathrm{rob}}\leq100M_i$。所有组值同为 $c$ 时，得分为 $100c$；任一组值上升而其余不变时，得分不会下降。为限制优势维度在五维中为薄弱维度过度补偿，从而无法在保留平均表现的情况下考虑最弱一项，我们采用以模型最弱维度限制优势维度的补偿，不修改原始指标的方式作为补偿规则，故而可以使连续评分避免样本刚跨过冲突阈值就突然扣分。取 $\lambda=0.25$ ，原 16 维 TOPSIS 的差别同时包含指标选择和聚合方式变化，只有相对于五维加权平均的式~\eqref{eq:q12-closed} 分差才考虑归于补偿规则。

\subsection{A1 与扩展集结果}

\begin{table}[htbp]
\centering
\caption{冲突覆盖情况}
\label{tab:q12-coverage}
\begin{tabular}{lrrr}
\toprule
分区 & 样本数 & 冲突数 & 冲突率（\%） \\
\midrule
A1 全量         & 51,230  & 15,124 & 29.5218 \\
A1 留出集       & 10,311  & 3,039  & 29.4734 \\
A2 新增         & 16,104  & 9,408  & 58.4203 \\
A3 新增         & 193,752 & 65,795 & 33.9584 \\
A1--A3 唯一 ID & 261,086 & 90,327 & 34.5966 \\
\bottomrule
\end{tabular}
\end{table}

A1 四类样本数依次为 15,124、14,958、14,439、6,709，合计 51,230。拟合集冲突率 $29.5340\%$，留出集 $29.4734\%$，二者接近只说明本次划分下判别比例相近。arxiv 域在 A1 和新增集的冲突率分别为 $56.24\%$ 和 $58.42\%$；github 域分别为 $33.47\%$ 和 $33.96\%$。A2 新增中“整洁度高、非重复程度低”占 $54.53\%$，A3 新增中反向组合占 $26.27\%$，表明主要冲突模式随领域而变。其余五域没有扩展集可作同检验。把全量域级冲突率按 A1 领域比例重加权后为 $29.6680\%$，接近 A1 的 $29.5218\%$；总体率升至 $34.5966\%$ 主要与领域构成变化相关。图~\ref{fig:q12-pairs} 展示有向冲突组合，图~\ref{fig:q12-domain} 展示 A1 各领域差异，为量化质量信号冲突在不同领域中的分布差异，采用 A1 拟合集确定的统一分位阈值，将五个质量维度中同时存在高分与低分信号的文本标记为冲突，并统计各领域的冲突率及样本量。采用领域内比例有助于比较不同规模领域的冲突发生频率，同时报告样本量可以明确比例对应的数据基础。结果表明，A1 总体冲突率为$29.52\%$，各领域存在较大的描述性差异：arXiv 为 $56.24\%$，高于其他领域；StackExchange 为 $15.35\%$，相对较低。Book 的冲突率为 $46.78\%$。C4 的冲突率虽低于 arXiv，但其样本量较大，冲突数达到 3,939 条，为各领域最多。总体体现为，高低质量信号并存、频繁程度随领域而异。

同一文本可能计入多个高低分配对，如图~\ref{fig:q12-pairs} 中单元格所示。
\begin{figure}[htbp]
\centering
\includegraphics[width=\linewidth]{fig3_1_conflict_pair_heatmap.pdf}
\caption{高低分信号配对比例}
\label{fig:q12-pairs}
\end{figure}

\begin{figure}[htbp]
\centering
\includegraphics[width=0.92\linewidth]{fig3_2_domain_conflict.pdf}
\caption{A1 各领域冲突率及样本量}
\label{fig:q12-domain}
\end{figure}

\begin{table}[htbp]
\centering
\caption{均值分对比（下降量为模型规定的保守折减）}
\label{tab:q12-scores}
\begin{tabular}{lrrrr}
\toprule
分区 & 原 TOPSIS & 加权平均 & 有限补偿分 & 相对基线变化 \\
\midrule
A1 全量         & 60.9471 & 58.6944 & 52.8212 & $-5.8733$ \\
A1 留出集       & 61.0370 & 58.7548 & 52.8643 & $-5.8905$ \\
A2 新增         & 63.4735 & 76.5370 & 72.2455 & $-4.2915$ \\
A3 新增         & 51.6887 & 46.9338 & 40.4316 & $-6.5022$ \\
A1--A3 唯一 ID & 54.2323 & 51.0674 & 44.8250 & $-6.2424$ \\
\bottomrule
\end{tabular}
\end{table}

A2 新增有限补偿分与原 TOPSIS 分的 Spearman 相关为 0.6105，说明换指标与换聚合规则会明显改变排序。若将候选分作为后续配比或标度律输入，原质量评分与补偿风存在一定差异可能需要分别讨论。

\subsection{灵敏度、核验与适用边界}

图~\ref{fig:q12-threshold} 与图~\ref{fig:q12-lambda} 分别展示两类敏感性。第一种是冲突率对分位阙值的敏感性，见图可知把尾部分位从 $10\%$ 调至 $25\%$，A1 留出集冲突率由 $9.4074\%$ 升至 $41.4606\%$，A2 新增由 $3.8934\%$ 升至 $79.0673\%$，A3 新增由 $13.7919\%$ 升至 $44.2380\%$，冲突率高度依赖相对阈值，$\lambda$ 在 $0,0.1,0.25,0.5$ 上的变化影响候选均分。第二种是$\lambda$ 对候选均分的影响，~\eqref{eq:q12-closed} 还给出解析灵敏度 $\partial Q_i^{\mathrm{rob}}/\partial\lambda=-100(M_i-m_i)$，维度差距越大，保守程度的影响越强；最弱维度若受噪声或域失配影响，模型也会随之偏低。

\begin{figure}[htbp]
\centering
\includegraphics[width=0.92\linewidth]{fig3_3_threshold_sensitivity.pdf}
\caption{冲突率对分位阈值的敏感性}
\label{fig:q12-threshold}
\end{figure}

\begin{figure}[htbp]
\centering
\includegraphics[width=0.92\linewidth]{fig3_4_compensation_sensitivity.pdf}
\caption{$\lambda$ 对均分的影响}
\label{fig:q12-lambda}
\end{figure}

原评分复算的七域均分最大误差小于 $10^{-12}$；重复 ID 的指标和领域一致，四类判别穷尽样本，式~\eqref{eq:q12-closed} 的闭式解与单调性可直接核对。


